\subsection{Fine-Tuning Datasets and Training Configuration}

We instantiate the agentic model $\Phi_w$ with
\texttt{deepseek-ai/DeepSeek-R1-Distill-Qwen-14B} and adapt it using LoRA. The
base-model parameters remain frozen during fine-tuning, while the trainable
parameters are restricted to the LoRA adapter weights. The checkpoint used in
the experiments is denoted by \texttt{checkpoint-850}.

The supervised fine-tuning corpus consists of four instruction-answer datasets.
Let
\[
    \mathcal{D}
    =
    \mathcal{D}_{\mathrm{state}}
    \cup
    \mathcal{D}_{\mathrm{action}}
    \cup
    \mathcal{D}_{\mathrm{incident}}
    \cup
    \mathcal{D}_{\mathrm{local}},
\]
where each element $(\mathbf{x}^i,\mathbf{y}^i)\in\mathcal{D}$ contains an
instruction prompt $\mathbf{x}^i$ and a target response $\mathbf{y}^i$.

\begin{itemize}
    \item \textbf{State prediction dataset}
    $\mathcal{D}_{\mathrm{state}}$. We use \texttt{states\_examples.json} from
    \texttt{kimhammar/CSLE-IncidentResponse-V1}. This dataset provides
    supervision for predicting incident-response state transitions.

    \item \textbf{Recovery action dataset}
    $\mathcal{D}_{\mathrm{action}}$. We use \texttt{action\_examples.json} from
    \texttt{kimhammar/CSLE-IncidentResponse-V1}. This dataset provides
    supervision for generating high-level response and recovery actions.

    \item \textbf{Incident understanding dataset}
    $\mathcal{D}_{\mathrm{incident}}$. We use \texttt{incident\_examples.json}
    from \texttt{kimhammar/CSLE-IncidentResponse-V1}. This dataset provides
    supervision for extracting incident descriptions, involved entities, and
    security-relevant labels from system and log evidence.

    \item \textbf{Local transformed incident classification dataset}
    $\mathcal{D}_{\mathrm{local}}$. We use a local JSON dataset,
    \texttt{transformed\_dataset\_cls\_pri\_all2preprocessing.json}, to provide
    additional supervision for identifying targeted systems and
    recovery-relevant assets.
\end{itemize}

In the final four-dataset training run, we sample at most 5,000 examples from
$\mathcal{D}_{\mathrm{state}}$, 5,000 examples from
$\mathcal{D}_{\mathrm{action}}$, 15,000 examples from
$\mathcal{D}_{\mathrm{incident}}$, and 5,000 examples from
$\mathcal{D}_{\mathrm{local}}$. This yields at most 30,000 supervised
instruction-answer pairs. The combined training set is shuffled using random
seed 99125.

The LoRA configuration is given by rank $r=64$, scaling factor
$\alpha=128$, and dropout probability $0.05$. The adapter is applied to the
\texttt{q\_proj} and \texttt{v\_proj} modules, with no trainable bias terms.
The resulting adapter contains 50,331,648 trainable parameters, corresponding
to approximately 0.36\% of the 14B-parameter base model.

The training uses fused AdamW with learning rate $9.5\times 10^{-4}$ and a
linear learning-rate schedule. The per-device batch size is 1, with 32 gradient
accumulation steps, giving an effective batch size of 32. We train for one
epoch using bfloat16 precision, maximum gradient norm 1.0, no warmup steps, and
no weight decay. We log the training loss at every optimization step and save a
checkpoint every 50 steps, retaining the latest three checkpoints. The
checkpoint used in the experiments is taken at global step 850, which
corresponds to approximately 0.907 training epochs, with logged training loss
0.6054.

